\section{Roles y permisos}

Nueve roles definidos en el DDL y en \codigo{Domain/Entidades/Seguridad.cs}.
Los permisos se guardan en \codigo{rol\_permiso} como pares
\codigo{AREA:ACCION} (leer, crear, editar, validar, administrar, exportar) y
se entregan como claims \codigo{permiso} vía
\codigo{PermisoPolicyProvider}; el menú solo muestra lo autorizado y el
servidor deniega y registra (\codigo{DENEGADO}) lo no autorizado.

\begin{longtable}{@{}p{0.22\textwidth}p{0.70\textwidth}@{}}
\toprule
\textbf{Rol} & \textbf{Permisos vigentes (pasos 03, 14 y 15)}\\
\midrule
SECRETARIA & Crea, edita y formaliza en GENERAL/FAMILIA/RETIRO/
  AUTORIZACIONES/DOCUMENTOS/SALUD/JUDICIAL; exporta. No valida.\\
DIRECCION & Lee todas las áreas; valida y exporta. No crea ni edita.\\
INSPECTORIA & Lee todas las áreas; edita datos del estudiante
  (GENERAL/FAMILIA/RETIRO). No crea ni valida.\\
EQUIPO\_PIE & Solo lee GENERAL/FAMILIA/PIE/DOCUMENTOS.\\
TENS & Lee GENERAL/FAMILIA/SALUD/MEDICACION/DOCUMENTOS; crea dosis
  (\codigo{MEDICACION:CREAR}).\\
UTP & Lee todas las áreas; exporta. Sin escritura.\\
CONVIVENCIA & Solo lee GENERAL/FAMILIA/RETIRO/AUTORIZACIONES/DOCUMENTOS
  (sin salud, PIE, judicial ni medicación).\\
DOCENTE & Solo \codigo{GENERAL:LEER}: datos mínimos (identificación, curso,
  estado).\\
ADMIN\_TI & \codigo{ADMIN:ADMINISTRAR/LEER} y \codigo{AUDITORIA:LEER};
  sin acceso al expediente.\\
\bottomrule
\end{longtable}

\pendiente{La matriz es la propuesta de Johan (17--18/11/2026) cargada en los
pasos 03, 14 y 15; la checklist la marca pendiente de confirmación final del
colegio. No modificar \codigo{rol\_permiso} hasta esa confirmación.}
